% CLAUDE REPORT::DIFFICULTY-LOCALIZATION-0001::2026-11-12T13:00Z
% Independent Claude contribution (not an LSPV role report; owns no main.tex region).
\documentclass[11pt]{article}
\usepackage{amsmath,amssymb,amsthm}
\usepackage[margin=1in]{geometry}
\usepackage{hyperref}
\newtheorem{theorem}{Theorem}
\newtheorem{lemma}[theorem]{Lemma}
\newtheorem{corollary}[theorem]{Corollary}
\theoremstyle{definition}
\newtheorem{definition}[theorem]{Definition}
\newcommand{\Wbar}{\overline W}
\title{Difficulty localization for the worst-case minimum-derivation program:\\
the two Route~B arms move in opposite directions along p-simulation}
\author{Claude (independent contributor), run DIFFICULTY-LOCALIZATION-0001}
\date{2026-11-12, 13:00 UTC (05:00 PST)}
\begin{document}
\maketitle

\noindent\textbf{Status summary.} Lemma~\ref{lem:mono}: PROVED/EXACT (elementary).
Theorem~\ref{thm:x3}: CONDITIONAL on the project result LSPV-1/375 ($X3\equiv_p ER_2$; specialist review still pending)
together with KNOWN THEOREMS (Cook--Reckhow 1979; standard ER/EF equivalence for CNF refutations).
Corollary~\ref{cor:popt}: KNOWN THEOREM repackaged (Cook--Reckhow; Krajíček--Pudlák 1989).
No lower bound and no bridge is proved. \textbf{P versus NP remains OPEN}; every load-bearing Route~B arrow
([P2-B1] superpolynomial $\Wbar_R$, [P2-B2]/B3 noncircular same-$R$ bridge) remains \textbf{OPEN}.

\section{Fixed objects and quantifiers}
Follow the 2026-11-02 refocus directive verbatim. For a fixed sound, complete, effectively specified
system $R$ with nonnegative node weights,
\[
W_R(F)=\min_{\pi:F\vdash_R\bot}\mathrm{cost}_R(\pi),\qquad
\Wbar_R(N)=\max_{F\in\mathrm{UNSAT},\ |\langle F\rangle|\le N} W_R(F).
\]
Write $\mathrm{LB}(R)$ for ``$\Wbar_R$ is superpolynomial'' (for all $C,k,N_0$ some $N\ge N_0$ has
$\Wbar_R(N)>C(N+1)^k$) and $\mathrm{B3}(R)$ for ``$\mathrm{SAT}\in\mathrm P\Rightarrow \exists C,k\,\forall N:\ \Wbar_R(N)\le C(N+1)^k$''.
Route~B for $R$ is the conjunction $\mathrm{LB}(R)\wedge\mathrm{B3}(R)$, which implies $\mathrm{SAT}\notin\mathrm P$.
Free optimal selection of the derivation is granted throughout; no proof-search cost is charged.

\begin{definition}
$R'$ \emph{cost-simulates} $R$, written $R\le_c R'$, if there is a polynomial $p$ with
$W_{R'}(F)\le p(|\langle F\rangle|+W_R(F))$ for every UNSAT CNF $F$. (Existence of short proofs only;
no efficient translation is required.)
\end{definition}

\section{The opposite-direction lemma}
\begin{lemma}[Monotonicity, PROVED/EXACT]\label{lem:mono}
If $R\le_c R'$ then (i) $\mathrm{LB}(R')\Rightarrow\mathrm{LB}(R)$ and (ii) $\mathrm{B3}(R)\Rightarrow\mathrm{B3}(R')$.
\end{lemma}
\begin{proof}
WLOG $p$ is nondecreasing. If $\Wbar_R(N)\le C(N+1)^k$ for all $N$, then for every UNSAT $F$ with
$|\langle F\rangle|\le N$, $W_{R'}(F)\le p(N+C(N+1)^k)$, so $\Wbar_{R'}$ is polynomially bounded. This is
the contrapositive of (i), and applied under the hypothesis $\mathrm{SAT}\in\mathrm P$ it gives (ii).
\end{proof}
\noindent\emph{Reading.} Strengthening the fixed calculus makes the bridge easier and the lower bound
harder; weakening does the reverse. A Route~B proof needs one calculus at which both arms hold. The
project's history matches this exactly: Haken/Urquhart give $\mathrm{LB}$ for ordinary Resolution, for which
$\mathrm{B3}(\mathrm{Res})\iff\mathrm{SAT}\notin\mathrm P$ (already recorded: LSPV-1/339, REFEREE331);
adding GF(2), AND3 and DEF ($R_{\mathrm{mix341}}$) moved the family failures (LSPV-2/341 Tseitin $4h-1$,
LSPV-3/343 bit-PHP polynomial via DEF and Cook's ER proofs).

\section{Where the lower-bound arm sits for the extension-strength calculus}
\begin{theorem}[CONDITIONAL on LSPV-1/375]\label{thm:x3}
Let $R$ be any fixed calculus with $ER_2\le_c R\le_c EF$, in particular $R=X3$ of LSPV-1/375 (and $R=ER_2$).
Then
\[
\mathrm{LB}(R)\iff \text{EF is not polynomially bounded on refutations of CNFs (measured in proof lines, hence also in size)}.
\]
\end{theorem}
\begin{proof}[Proof sketch]
KNOWN: ER with arbitrary-width definitions and $ER_2$ (two-input definitions) are polynomially
equivalent (split a wide conjunction into a chain of binary definitions); ER and EF are polynomially
equivalent on CNF refutations (Cook--Reckhow 1979; Krajíček, \emph{Bounded Arithmetic, Propositional Logic and
Complexity Theory}, 1995, \S4.5). Lines versus bit size: every ER line mentions at most $N+q$ variables, so a
$q$-line ER refutation has size $\le\mathrm{poly}(N+q)$ (this is the same bookkeeping as LSPV-4/342 for
$R_{\mathrm{mix341}}$). Hence $ER_2\equiv_c EF$ in the line measure, and the sandwich plus
Lemma~\ref{lem:mono}(i) in both directions gives the equivalence. The upper inclusion $X3\le_c ER_2$ is
the content of LSPV-1/375, still under specialist review; that is the only non-classical premise.
\end{proof}

\begin{corollary}[consequences for planning]\label{cor:x3}
For $R\in\{ER_2,X3\}$, [P2-B1] is \emph{exactly} the problem of a superpolynomial Extended Frege lower bound for
some CNF family. No superpolynomial lower bound is known even for $AC^0[p]$-Frege for any prime $p$, let alone
EF. Interpolation-based routes are known to fail at EF strength under standard cryptographic assumptions
(Krajíček--Pudlák 1998, RSA; Bonet--Pitassi--Raz 2000 for Frege, Diffie--Hellman). So in this window the
lower-bound arm is not a ``weaker'' substitute for P$\ne$NP-type difficulty; it is the central open problem of
proof complexity. Conversely, its bridge $\mathrm{B3}(EF)$ is classically sufficient-conditioned: it follows
if SAT$\in$P via an algorithm whose correctness is provable in $PV$/$S^1_2$ (Cook 1975; already recorded as
LSPV-1/357's firewall, which also shows provability, not truth, is what is used).
\end{corollary}

\begin{corollary}[KNOWN, repackaged: the top of the order]\label{cor:popt}
If a p-optimal (or merely optimal) proof system $R^\ast$ exists, then $\mathrm{B3}(R^\ast)$ holds
unconditionally and $\mathrm{LB}(R^\ast)\iff\mathrm{NP}\ne\mathrm{coNP}$.
\end{corollary}
\begin{proof}
SAT$\in$P implies NP$=$coNP, hence some proof system is polynomially bounded (Cook--Reckhow), hence
the optimal $R^\ast$ is (simulating it). The second claim is Cook--Reckhow again. Existence of optimal systems is
OPEN (Krajíček--Pudlák 1989: implied by NE$=$coNE). Note P$=$NP implies E$=$NE and hence a p-optimal system,
but that system is constructed from the hypothetical algorithm; it is not a preselected $R$, so the
directive's warning stands.
\end{proof}

\noindent\textbf{Net localization.} Along the simulation order the program has three regimes:
(a) below EF strength where lower bounds are known (Res, tree-like/regular/bounded-depth Res($\oplus$),
$R_{\mathrm{mix341}}$ minus DEF on restricted families): $\mathrm{LB}$ known or near, $\mathrm{B3}\iff$P$\ne$NP;
(b) at ER/EF strength ($ER_2$, $X3$, and $R_{\mathrm{mix341}}$ if its full LOCAL-CONVERT is normalized):
$\mathrm{LB}\iff$ EF lower bound (open), $\mathrm{B3}\Leftarrow$ PV-provable SAT$\in$P;
(c) at the top: $\mathrm{B3}$ free, $\mathrm{LB}\iff$NP$\ne$coNP. No regime gives a ``cheap'' arm.
This does not refute the program; it says which known open problem each arm is.

\section{Falsifiers proposed}
\begin{enumerate}
\item[F1] A calculus $R$ with $\mathrm{B3}(R)$ provable unconditionally and a known $\mathrm{LB}(R)$ would
contradict nothing here but would prove P$\ne$NP; any claim of both must be checked against
Lemma~\ref{lem:mono} for a hidden change of $R$ between the two arms (the standard failure mode).
\item[F2] If LSPV-1/375's compiler fails (an X3 shortcut not polynomially compilable into $ER_2$), the
upper half of Theorem~\ref{thm:x3} fails for $X3$ and $\mathrm{LB}(X3)$ could be strictly harder than EF.
\item[F3] Full $R_{\mathrm{mix341}}$ with opaque circuit-valued LOCAL-CONVERT payloads may exceed EF strength
(LSPV-1/375 F3, LSPV-4/363); then it sits between regimes (b) and (c) and the theorem does not apply.
\end{enumerate}

\section{Literature checked this run}
Queries: ``2026 proof complexity lower bound Res($\oplus$)/extended Frege/bounded-depth Frege'';
``Pich Santhanam Towards P$\ne$NP from Extended Frege lower bounds''. Findings:
Itsykson, Podolskii, Shekhovtsov, ECCC TR26-018 / CCC 2026 (LIPIcs 383, paper 13): simplified and extended
lifting for bounded-depth Res($\oplus$) and a pure quadratic Res($\oplus$) size lower bound; unrestricted
superpolynomial Res($\oplus$) lower bounds remain open (regime (a), does not reach $R_{\mathrm{mix341}}$; already
cited once in the corpus). Pich--Santhanam, arXiv:2312.08163: EF lower bounds for circuit-lower-bound
tautologies plus their hypotheses I--II imply P$\ne$NP; Arteche--Khaniki--Pich--Santhanam (ICALP 2024) note this
would give NEXP$\not\subseteq$P/poly. Both are already in the corpus (LSPV-1/416, codex 2026-11-11); they are the
closest published analogue of a regime-(b) Route~B and confirm its difficulty. Alekseev, Theory of Computing 22
(2026): polynomial calculus with extension simulating EF (regime (b), no lower bound). No 2026 EF lower bound found.
Barrier-transfer classification: interpolation $\to$ FAILED TRANSFER at regime (b) under RSA/DH; bounded-arithmetic
reflection $\to$ DIRECTLY TRANSFERABLE THEOREM for the conditional B3 sufficient condition only.

\section{Corpus reviewed; duplication audit}
Read: root \texttt{AGENTS.md}; the 2026-11-02 worst-case refocus directive with all embedded results
(LSPV-2/341, /342; LSPV-3/342, /343; LSPV-4/341, /342; REFEREE295); the 2026-11-09 knot/link handoff; the five
2026-11-12 handoffs (LSPV-1/589, LSPV-2/573, LSPV-3/587, LSPV-4/559, LSPV-5/502); LSPV-1/339, /357, /375;
REFEREE331; root \texttt{main.tex} status lines (2.19\,MB; grepped, not reread in full).
\textbf{Duplicates:} ``B3$\iff$SAT$\notin$P for a lower-bounded $R$'' is already in LSPV-1/339 and REFEREE331;
the reflection sufficient condition is LSPV-1/357; $X3\equiv_p ER_2$ is LSPV-1/375. \textbf{New here:} the
two-sided monotonicity lemma stated as a planning tool, the explicit identification $\mathrm{LB}(X3)\iff$ EF
lower bound, and the three-regime map. These are elementary and may already be implicit in peer text; no
novelty beyond that is claimed. \textbf{Stale/overstated claims found:} none. \textbf{Root main.tex:} not edited
(no Claude-owned region); this report is evidence for LSPV-5 reconciliation only. The knot/link handoff:
TRACKERS REVIEWED; NO MATERIAL CHANGE from this run.

\noindent P versus NP remains OPEN.
\end{document}
% END_OF_VERIFIED_REPORT::CLAUDE::DIFFICULTY-LOCALIZATION-0001
